\documentclass[11pt,a4paper]{article}
\usepackage[utf8]{inputenc}
\usepackage[T1]{fontenc}
\usepackage{lmodern,amsmath,amssymb,textcomp}
\usepackage[margin=24mm]{geometry}
\usepackage{longtable,booktabs,array,enumitem}
\usepackage[hidelinks]{hyperref}
\setlength{\parindent}{0pt}
\setlength{\parskip}{6pt}
\emergencystretch=3em
\hypersetup{pdfauthor={Rachel Teo, Wenyi Wang, Hamdi Abderrahmene, Alejandro Zarzuelo Urdiales}}
\begin{document}
{\LARGE\bfseries Exact admissible powers and a sharp logarithmic boundary scale for a Hermitian quartic\par}\medskip
{\small Rachel Teo\ensuremath{^1}, Wenyi Wang\ensuremath{^1}, Hamdi Abderrahmene\ensuremath{^1}, Alejandro Zarzuelo Urdiales\ensuremath{^2}\par}\medskip
{\small\textbf{Affiliations:} \href{https://sakana.ai/}{\ensuremath{^1} Sakana AI}; \href{https://rsihouse.ai/openmath}{\ensuremath{^2} Open Math}.\par}

{\small\href{https://github.com/alejandrozu/openmath-2026-judging}{Code and formal proofs}\par}

\subsection{Abstract}
For H\_t(z)=(1+|z|\ensuremath{^2})\textasciicircum{}4\ensuremath{-}t|z|\ensuremath{^4}, the positive exponents producing a holomorphic squared-norm representation are classified exactly on6<t\ensuremath{\leq}7. Near t=8, the least positive admissible exponent has scale log(1/\ensuremath{\delta})/\ensuremath{\delta} with leading constant one, where \ensuremath{\delta}=8\ensuremath{-}t.

\subsection{Squared norms and the finite exponent classification}
A function-level squared norm means a finite sum \ensuremath{\sum}|f\_j(z)|\ensuremath{^2} with holomorphic polynomials f\_j, without requiring the summands to be monomials. For a radial polynomial in |z|\ensuremath{^2}, the formal development proves equivalence with nonnegativity of every coefficient in the one-variable polynomial. The necessity is a Gram/coefficient statement, so the result concerns actual functions rather than only a chosen monomial ansatz.

For 6<t\ensuremath{\leq}7 and a positive integer N, H\_t\textasciicircum{}N has such a representation exactly when N=2, or N\ensuremath{\geq}4, or N=3 and 924\ensuremath{-}210t+18t\ensuremath{^2}\ensuremath{-}t\ensuremath{^3}\ensuremath{\geq}0. The first power fails because its middle coefficient 6\ensuremath{-}t is negative. Exact expansions show that the square and fifth power have nonnegative coefficients throughout this interval. Products of those powers cover every exponent at least four. In the cubic expansion the only possible obstruction is the displayed central coefficient.

\subsection{Lower bounds near the radical boundary}
For 6<t<8, an admissible exponent satisfies 12N(8\ensuremath{-}t)\ensuremath{\geq}t(10\ensuremath{-}t). This inverse-distance bound comes from signed coefficient moments and positivity constraints. It is strengthened near t=8 by a Fourier obstruction: a negative terminal interval in the trigonometric profile must be compensated by the positive majorant allowed by nonnegative coefficients.

The source derives 3t\textasciicircum{}N N\textasciicircum{}(1/4)\ensuremath{\leq}28(16\ensuremath{-}t)\textasciicircum{}N, equivalently 81N t\textasciicircum{}(4N)\ensuremath{\leq}614656(16\ensuremath{-}t)\textasciicircum{}(4N), for 7\ensuremath{\leq}t<8. Taking logarithms and combining with the elementary lower bound gives the quantified leading lower scale. For every \ensuremath{\varepsilon}>0, sufficiently small \ensuremath{\delta}=8\ensuremath{-}t>0 forces \ensuremath{\delta}N\ensuremath{\geq}(1\ensuremath{-}\ensuremath{\varepsilon})log(1/\ensuremath{\delta}) for every positive admissible N.

\subsection{Matching upper bound and least exponent}
The upper bound extracts each coefficient by Fourier integration at a positive radius selected to match that coefficient's mean. The source proves existence of the saddle radius, the local Gaussian expansion of the actual normalized quartic, the modulus gap on outer arcs, and bounds for Gaussian, quartic and edge errors. Small coefficients receive a separate direct finite estimate, and symmetry covers the opposite edge. These steps establish positivity in every degree 0,...,4N uniformly whenever N\ensuremath{\geq}(1+\ensuremath{\varepsilon})log(1/\ensuremath{\delta})/\ensuremath{\delta} and \ensuremath{\delta} is sufficiently small.

Define the least positive admissible index using that nonempty admissible set. The formal squeeze proves both its minimality and 1\ensuremath{-}\ensuremath{\varepsilon}\ensuremath{\leq}\ensuremath{\delta}\ensuremath{\cdot}index/log(1/\ensuremath{\delta})\ensuremath{\leq}1+\ensuremath{\varepsilon} near zero. It excludes exponent zero, whose constant polynomial would otherwise trivialize the minimum. The claimed endpoint is the quantified squeeze, rather than a separately formalized filter-limit statement or eventual-index theorem.

\subsection{Prior work and scope}
The quartic family, threshold t<8, t=7 square and qualitative eventual positivity are known background. The precise refinements considered here are the exceptional exponent/cubic transition and logarithmic boundary scale with leading constant one. The inspected Tan-To appendix concerns an auxiliary power R\textasciicircum{}m and a different \ensuremath{\Theta}(1/\ensuremath{\alpha}) scale, which does not itself imply this H\_t\textasciicircum{}N statement. Comparison with its full main theorem and the Handelman/To-Yeung literature remains incomplete.

The accompanying analytic exposition derives variance certificates, the local complex logarithm, positive Gaussian core, modulus deficit, small-radius estimates and finite-window assembly for the displayed one-variable theorem. No additional infinite analytic lemma is an unstated premise of that result. The general Hermitian-power problem and a formally verified binary homogeneous bridge remain outside its scope.

\subsection{Formal verification}
The 40-source development was checked with Lean; the selected exact-exponent and boundary statements use standard kernel axioms. Exact compiler/library pins, theorem sources and audit records are supplied in the companion repository. This verification does not settle the separate historical-priority comparison.

\begin{samepage}
\subsection{PowerSemigroup.function\_exact\_exponent\_semigroup}
\begin{quote}\small\ttfamily theorem function\_exact\_exponent\_semigroup \{t : \ensuremath{\mathbb{R}}\} (ht6 : 6 < t) (ht7 : t \ensuremath{\leq} 7)\newline     (n : \ensuremath{\mathbb{N}}) (hn : 0 < n) :\newline     FunctionSquaredNorm (fun z : \ensuremath{\mathbb{C}} =>\newline       (((1+Complex.normSq z)\textasciicircum{}4-t*(Complex.normSq z)\textasciicircum{}2)\textasciicircum{}n : \ensuremath{\mathbb{R}})) \ensuremath{\leftrightarrow}\newline       n=2 \ensuremath{\vee} 4\ensuremath{\leq}n \ensuremath{\vee} (n=3 \ensuremath{\wedge} 0\ensuremath{\leq}924-210*t+18*t\textasciicircum{}2-t\textasciicircum{}3)\end{quote}
{\small Source: proofs/opdp13/OpHack/PowerSemigroup/Evaluation.lean:116.\par}

{\small Source SHA-256: 9924251f0abcd263e38621ad0b96cadbd0a33f6c353d41ecf3e16e30a860ebfc\par}

\end{samepage}
\begin{samepage}
\subsection{PowerMoment.function\_exponent\_lower\_bound}
\begin{quote}\small\ttfamily theorem function\_exponent\_lower\_bound (t : \ensuremath{\mathbb{R}}) (ht : 6 < t) (ht8 : t < 8)\newline     (n : \ensuremath{\mathbb{N}}) (hn : 0 < n)\newline     (h : PowerSemigroup.FunctionSquaredNorm (fun z : \ensuremath{\mathbb{C}} =>\newline       (((1+Complex.normSq z)\textasciicircum{}4-t*(Complex.normSq z)\textasciicircum{}2)\textasciicircum{}n : \ensuremath{\mathbb{R}}))) :\newline     t*(10-t) \ensuremath{\leq} 12*(n : \ensuremath{\mathbb{R}})*(8-t)\end{quote}
{\small Source: proofs/opdp13/OpHack/PowerMoment/Bound.lean:86.\par}

{\small Source SHA-256: 556b24cda4326f34d2b8e723b629b76a7930ff6a4c115f831942550fdd5fa830\par}

{\small\textbf{Sources and attribution:} \href{https://aimath.org/WWN/complexitycrmap/complexitycrmap.pdf}{1. aimath.org / complexitycrmap.pdf}; \href{https://arxiv.org/abs/1012.2479}{2. arXiv source: 1012.2479}; \href{https://doi.org/10.1007/s12220-019-00334-9}{3. DOI: 10.1007/s12220-019-00334-9}; \href{https://doi.org/10.1007/s12220-025-02228-5}{4. DOI: 10.1007/s12220-025-02228-5}; \href{https://github.com/alejandrozu/openmath-2026-judging/blob/d0ed487f487fa7ec814ff705ab55249983fe8707/OpenMath-Judging/review/entries/E02-opdp13-quartic.txt}{5. alejandrozu/openmath-2026-judging / E02-opdp13-quartic.txt}; \href{https://github.com/alejandrozu/openmath-2026-judging}{Proof source repository}.\par}

\end{samepage}
{\small \textbf{AI assistance.} Codex assisted literature checking, editorial preparation and analysis of the supplied formal artifacts. The human authors are responsible for checking the final mathematical claims, references and attribution.\par}

\end{document}
